\section{Discussion}

RevTrack suggests that review assistance should be evaluated as evidence tracking rather than static critique generation.
The practical failure mode is not merely that a model makes a wrong label.
It is that the model may keep repeating a criticism that has been fixed, or may credit a response that only defers the issue.
Both failures can distort author effort and reviewer attention.

The benchmark also changes how model progress should be read.
Small accuracy differences are not meaningful when most examples are partially fixed.
Macro-F1 and per-label recovery are closer to the scientific workload: fixed cases test whether stale criticism is retired, while unresolved cases test whether a model resists superficial reassurance.
The failure taxonomy makes this point concrete.
On the expanded ICLR 2025 frontier, many transfer errors are not random confusions; they are recognizable scientific-assistance failures such as over-crediting unresolved concerns, keeping fixed criticisms alive, and missing regressions.

The expanded ICLR 2025 frontier strengthens the cross-year story but also sharpens the boundary of the claim.
It confirms brittleness on a larger active frontier, yet the sample is deliberately hard and disagreement-heavy.
This is useful evidence for a benchmark paper because it shows where static matching breaks under pressure.
The NeurIPS active frontier and ICLR 2023 random/stratified slice broaden that evidence, but they still do not imply that the same label distribution holds across all venues or years.
A bounded independent second-pass mini60 check now provides reliability support, but the next step is still a non-ICLR random/stratified slice for broader prevalence-sensitive evidence.

More broadly, RevTrack points to a design requirement for scientific assistants.
They need an issue ledger: a structured representation of what was criticized, what evidence was added, and what status each concern should now have.
That representation is not a complete solution, but it gives the field a measurable target for longitudinal scientific judgment.
